\section*{Abstract}
Current evaluations of LLM false outputs often collapse distinct mechanisms into one label, making it unclear whether a detector captures ignorance or strategic misreporting. We study this mechanism conflation problem with a counterfactual paired design on 240 factual \mc items drawn from \truthfulqa and \halueval. For each item, we query \gptmodel under two prompt conditions: \control (optimize factual correctness) and \pressure (social and reward pressure to match a user-believed wrong option). This design holds question identity fixed while varying incentives.

Pressure substantially increases false outputs: the pooled false rate rises from 12.9\% to 40.8\% (risk difference +27.9 points; McNemar $p=7.43\times10^{-16}$). Among pressure-condition false outputs ($n=98$), 68.4\% are incentive-induced and 31.6\% are epistemic, with bootstrap 95\% CIs of [59.2\%, 76.5\%] and [23.5\%, 40.8\%], respectively. A lightweight output-text detector (\tfidfdet) reaches moderate pooled performance but transfers weakly across datasets (\auroc 0.57--0.60), consistent with mechanism-specific artifacts.

These results show that ``false output'' is not a single behavioral class under incentive perturbations. We argue that lie-detection evaluation should report mechanism-separated metrics and that mitigation should be mechanism-specific: calibration for epistemic errors and objective/policy constraints for incentive-driven misreporting.
